% Generated draft: requires mathematical and source review.
\documentclass[11pt]{article}
\usepackage[margin=25mm]{geometry}
\usepackage{fontspec}
\setmainfont{Latin Modern Roman}
\usepackage{amsmath,amssymb,mathtools,unicode-math}
\setmathfont{Latin Modern Math}
\usepackage{enumitem,needspace,xurl,hyperref}
\setlength{\parindent}{0pt}
\setlength{\parskip}{5pt}
\setlength{\emergencystretch}{3em}
\newcommand{\N}{\mathbb{N}}
\newcommand{\Z}{\mathbb{Z}}
\newcommand{\Q}{\mathbb{Q}}
\newcommand{\R}{\mathbb{R}}
\newcommand{\C}{\mathbb{C}}
\newcommand{\F}{\mathbb{F}}
\newcommand{\A}{\mathbb{A}}
\newcommand{\PP}{\mathbb{P}}
\newcommand{\E}{\mathbb{E}}
\newcommand{\eps}{\varepsilon}
\newcommand{\dd}{\,\mathrm{d}}
\begin{document}
\section*{UnsolvedMath ID: 30006340}
\textbf{Source problem code:} catalogue entry
\subsection{Definitions and mathematical statement}
For an integer \(n\geq 1\), \(S_n\) denotes the group of permutations of \(\{1,\ldots,n\}\), and \(A_n\) its subgroup of even permutations. A subgroup \(G\leq S_n\) is transitive if, for every \(i,j\in\{1,\ldots,n\}\), some \(g\in G\) satisfies \(g(i)=j\). A group is solvable if its derived series terminates at the trivial group. For \(x\in G\), its conjugacy class is \(\{gxg^{-1}:g\in G\}\); for subsets \(C_1,C_2,C_3\subseteq G\), their product is \(C_1C_2C_3=\{c_1c_2c_3:c_i\in C_i\}\).

No precise standalone mathematical problem can be recovered from the supplied record. Its supplied statement status is \texttt{unrecoverable}, and its clean statement is absent. The extraction is workshop-summary prose containing several different mathematical topics, without identifying one as the intended problem.

The opening fragment says that every transitive solvable subgroup of \(S_n\) was shown to have polynomial diameter, but it neither defines the diameter nor supplies the generating-set convention, polynomial parameter, or uniformity of the bound. Its introductory word ``conjecture'' does not identify a complete conjecture. Consequently this fragment cannot be promoted here to a quantified theorem or conjecture.

Another sentence reports that Daniele Dona discussed for which \(n\) the alternating group \(A_n\) can be covered by a product of three conjugacy classes. Although the notation above explains such a product, the source does not specify the intended range of \(n\), restrictions on the classes, known cases, or whether this sentence is the target of this registry entry. No separate problem is inferred from it.

The supplied verification note describes the extraction as workshop-summary prose ending at a dangling occurrence of ``conjecture'', whereas the supplied passage begins with that word and continues through workshop acknowledgements. This textual discrepancy remains unresolved. The registry's supplied label \texttt{open} comes from its unresolved-claim classification rule; it does not establish the existence or open status of a particular mathematical conjecture.
\subsection{Short English statement}
This entry preserves an unrecoverable workshop narrative. It mentions permutation-group diameter, quaternion-group presentations, orbit enumeration, products of conjugacy classes, subnormalisers, and generating functions, but supplies no uniquely identifiable mathematical problem. The supplied \texttt{open} label is a dated curation classification, and a precise statement requires human review of the intended source passage.
\subsection{Sources}
\begin{itemize}
\item Local UnsolvedMath registry, permanent ID \(30006340\), title ``Computational Group Theory Workshop Narrative'', exact source problem code \texttt{catalogue entry}; collection: Oberwolfach Reports Open Problems. The supplied record lists proposed year \(2025\), supplies no proposer, and retains the original passage for provenance.
\item Supplied report citation: \emph{Computational Group Theory} (2025), \emph{Oberwolfach Reports}, pp.\ 1389--1450, DOI \(10.4171/owr/2025/27\), \url{https://doi.org/10.4171/owr/2025/27}. The literature-sources entry labelled ``OWR report'' supplies this same link. This report citation is supplied provenance, not a source independently visited or verified for this draft; no precise problem location or authoritative standalone problem statement is supplied.
\item Literature assessment dated \texttt{2026-08-22}: the supplied curation labels the entry open under an unresolved-claim rule because it regards the extraction as a table-of-contents fragment followed by general workshop prose with no precise mathematical claim. The supplied background says this triage derives from OWR curation and calls for independent verification of time-sensitive status and cited claims. This is supplied dated curation, not an independently verified literature assessment.
\item The supplied statement-verification note reports that the official report confirms workshop-summary prose rather than a standalone problem, and that no clean statement was promoted. That verification is reported as source metadata, not independently confirmed here. Its description of the location of the dangling word ``conjecture'' differs from the supplied passage.
\item Attribution within the supplied narrative: Tommy Hofmann discussed Wall's \(D_2\) problem, a new family of quaternion-group presentations, and computational investigation of homotopy equivalence; Tobias Rossmann surveyed symbolic orbit enumeration for a group scheme acting on a scheme; Daniele Dona discussed products of three conjugacy classes in alternating groups; Gunter Malle reported properties of subnormalisers of \(p\)-elements connected with a recent character-theoretic conjecture attributed to Moretó and Rizo; Christopher Voll discussed multivariate generating functions for enumeration in algebra, geometry, and number theory. No precise conjecture or theorem citation accompanies these summaries.
\item Additional retained context: the narrative credits Alexander Hulpke's hybrid group framework, Chris Liu's matrix system solver, and Max Horn's OSCAR presentation during a two-hour computational-tools session. It mentions a problem session and collaborative discussions without stating their problems. The acknowledgement thanks the National Science Foundation for support of junior participants through grant \texttt{DMS-2230648}, ``US Junior Oberwolfach Fellows'', on behalf of the MFO and workshop organizers.
\end{itemize}
\end{document}
